%! AP Statistics — Unit 1, Lesson 1.3 — Guided Notes ANSWER KEY
\documentclass[10pt]{article}
\usepackage{apstats-article}
\usepackage{apstats-key}

%=============================================================================
\begin{document}

\pageheader{Unit 1, Lesson 1.3}{Guided Notes --- Answer Key}
\begin{tabularx}{\linewidth}{X X X}
Name:\;\ans{Key} & Date:\; & Period:\;
\end{tabularx}

\vspace{0.12in}

\begin{objectivebox}
\small By the end of this lesson, I will be able to\dots\\[4pt]
\ansline{construct a frequency and relative frequency table from raw categorical data.}\\
\ansline{interpret relative frequencies in context using complete sentences.}\\
\ansline{explain why relative frequency is required to compare groups of different sizes.}
\end{objectivebox}

\vspace{0.1in}

\begin{vocabbox}
\small
\noindent\textbf{\textcolor{navy}{Frequency:}}\\
\ans{The count of individuals in a particular category}\\[6pt]

\noindent\textbf{\textcolor{navy}{Relative frequency:}}\\
\ans{The frequency divided by the total number of individuals; expressed as a proportion
(decimal) or percent}\\[6pt]

\noindent\textbf{\textcolor{navy}{Frequency table:}}\\
\ans{A table listing each category of a categorical variable and the count of individuals
in that category}\\[6pt]

\noindent\textbf{\textcolor{navy}{Relative frequency table:}}\\
\ans{A table listing each category and the proportion (or percent) of individuals in that
category}\\[6pt]

\noindent\textbf{\textcolor{navy}{Distribution (of a categorical variable):}}\\
\ans{Describes all possible values a categorical variable takes and how often each value
occurs (given as relative frequencies)}\\[6pt]
\end{vocabbox}

\vspace{0.1in}

\begin{hookbox}
\small\textbf{Which school has more walkers?}

\vspace{8pt}
\renewcommand{\arraystretch}{1.6}
\begin{tabularx}{\linewidth}{@{} >{\bfseries}p{0.15\linewidth}
    C{0.12\linewidth} C{0.2\linewidth}
    C{0.12\linewidth} C{0.2\linewidth} @{}}
\rowcolor{sky}
\textbf{Mode} & \textbf{School A} & \textbf{Rel.\ Freq.\ A}
              & \textbf{School B} & \textbf{Rel.\ Freq.\ B} \\
\midrule
Bus  & 120 & \ans{0.50 (50\%)} & 30  & \ans{0.15 (15\%)} \\
Car  & 80  & \ans{0.33 (33\%)} & 50  & \ans{0.25 (25\%)} \\
Walk & 40  & \ans{0.17 (17\%)} & 120 & \ans{0.60 (60\%)} \\
\midrule
\rowcolor{sky}\textbf{Total} & 240 & \ans{1.00} & 200 & \ans{1.00} \\
\end{tabularx}

\vspace{10pt}
Which school has more walkers proportionally?\quad
\ans{School B (60\% vs.\ 17\%)}\\[6pt]
What does this tell you?\quad
\ans{Counts alone are misleading when groups have different sizes.
School A has more walkers by count (40 vs.\ 30) but far fewer by proportion.}

\begin{teachernote}
Emphasize: School A has 40 walkers but 240 total; School B has 120 walkers out of 200.
This is the core motivation for relative frequency.
\end{teachernote}
\end{hookbox}

\vspace{0.1in}

\begin{notesbox}{Building a Frequency and Relative Frequency Table}
\small
\begin{multicols}{2}

\textbf{\textcolor{navy}{Step-by-Step Process}}\\[4pt]
\begin{enumerate}[leftmargin=1.4em, itemsep=3pt]
  \item List all \ans{categories} in the first column.
  \item \ans{Tally} the data.
  \item Record the \ans{frequency} (count).
  \item Divide each count by \ans{$n$ (total)} to get relative frequency.
  \item \textbf{Check:} relative frequencies must sum to \ans{1}.
\end{enumerate}

\columnbreak

\textbf{\textcolor{navy}{Formula}}\\[4pt]
\[
  \text{Relative frequency} = \frac{\text{frequency}}{\ans{n \text{ (total)}}}
\]
\vspace{4pt}
Can be expressed as a \ans{proportion} (decimal)\\
or a \ans{percent} (multiply by 100).\\[8pt]
\textbf{\textcolor{navy}{Quick error check:}}\\
Sum of all relative frequencies $=$ \ans{1}\\
Sum of all relative frequencies $=$ \ans{100}\%

\end{multicols}
\end{notesbox}

\vspace{0.1in}

\begin{notesbox}{Worked Example --- Pet Ownership Survey ($n = 40$)}
\small

\renewcommand{\arraystretch}{1.8}
\begin{tabularx}{\linewidth}{@{} >{\bfseries}p{0.22\linewidth}
                                    C{0.16\linewidth}
                                    C{0.22\linewidth}
                                    C{0.22\linewidth} @{}}
\rowcolor{sky}
\textbf{Pet type} & \textbf{Tally} & \textbf{Frequency} & \textbf{Rel.\ Frequency} \\
\midrule
Dog    & \small||||||||||  & 12 & \ans{0.30 (30\%)} \\
Cat    & \small|||||||     & 8  & \ans{0.20 (20\%)} \\
Fish   & \small||||        & 6  & \ans{0.15 (15\%)} \\
Bird   & \small|||         & 4  & \ans{0.10 (10\%)} \\
None   & \small||||||||||  & 10 & \ans{0.25 (25\%)} \\
\midrule
\rowcolor{sky}\textbf{Total} & & \ans{40} & \ans{1.00} \\
\end{tabularx}

\vspace{10pt}
\textbf{AP-style sentence:}\\[4pt]
\ans{Of the 40 households surveyed, the most common pet type was dog, owned by
30\% of households. The least common was bird, owned by only 10\% of households.}

\begin{teachernote}
Accept any grammatically correct sentence that includes the total ($n = 40$),
the variable (pet type), the category, and the relative frequency as a percent.
\end{teachernote}
\end{notesbox}

\newpage

\begin{notesbox}{Describing a Distribution --- AP Language}
\small
\begin{multicols}{2}

\textbf{\textcolor{navy}{What to include:}}
\begin{itemize}[itemsep=3pt]
  \item The \ans{category} with the highest relative frequency.
  \item The \ans{category} with the lowest relative frequency.
  \item The \ans{total $n$} of the sample.
  \item \ans{Context} --- always include this.
\end{itemize}

\columnbreak

\textbf{\textcolor{navy}{Model sentence:}}\\[4pt]
\textit{``Of the \ans{40} households surveyed, the most common
pet was \ans{dog} with \ans{30}\% of responses.
The least common was \ans{bird} with \ans{10}\%.''}

\end{multicols}

\vspace{6pt}
\textbf{Distribution description (pet example):}\\[4pt]
\ans{Of the 40 households surveyed, dog was the most common pet type (30\%).
The second most common was ``none'' (25\%), followed by cat (20\%), fish (15\%),
and bird (10\% --- the least common).}

\end{notesbox}

\vspace{0.1in}

\begin{practicebox}
\small
\textbf{Transportation survey ($n = 25$):}

\vspace{8pt}
\renewcommand{\arraystretch}{1.8}
\begin{tabularx}{\linewidth}{@{} >{\bfseries}p{0.22\linewidth}
                                    C{0.18\linewidth}
                                    C{0.22\linewidth} @{}}
\rowcolor{sky}
\textbf{Mode} & \textbf{Frequency} & \textbf{Rel.\ Frequency} \\
\midrule
Car  & 10 & \ans{0.40 (40\%)} \\
Bus  & 8  & \ans{0.32 (32\%)} \\
Walk & 4  & \ans{0.16 (16\%)} \\
Bike & 3  & \ans{0.12 (12\%)} \\
\midrule
\rowcolor{sky}\textbf{Total} & \ans{25} & \ans{1.00} \\
\end{tabularx}

\vspace{10pt}
\begin{enumerate}[leftmargin=1.5em, itemsep=6pt]
  \item Proportion who walk: \ans{$4/25 = 0.16$}
  \item Percentage \textit{not} by bus: \ans{$1 - 0.32 = 0.68 = 68\%$}
        (or $(10+4+3)/25 = 17/25$)
  \item \ans{Of the 25 AP Statistics students surveyed, the most common mode of
        transportation was car, used by 40\% of students.}
\end{enumerate}

\end{practicebox}

\vspace{0.1in}

\begin{spiralbox}
\small
\begin{multicols}{2}

\textbf{This lesson connects forward to\dots}
\begin{itemize}[itemsep=2pt]
  \item \textbf{Lesson 1.4:} Bar charts and pie charts display the same distributions visually.
  \item \textbf{Unit 2:} Two-way tables extend this to two categorical variables.
  \item \textbf{Unit 4:} Relative frequency $\approx$ probability.
\end{itemize}

\columnbreak

\textbf{Big Idea:}\\[4pt]
\ans{When two groups have different total sizes, comparing raw counts is unfair ---
a group with 40 out of 50 walkers is very different from 40 out of 400. Relative
frequency puts both groups on the same scale (out of 1 or 100\%) so comparisons
are valid.}

\end{multicols}
\end{spiralbox}

\end{document}
